\documentclass[10pt,a4paper]{amsart}
\usepackage{amsmath,amssymb,amsthm,graphicx}
% Retained minimal preamble for the Stacks Project chapter more-algebra.tex (More on Algebra).
% The chapter file pinned-original-source.tex carries only \input{preamble} on its first line, so no
% preamble text lives inside the chapter itself.  The Stacks Project uses one single shared preamble for
% every chapter; the snapshot stages it identically as sources/stacks-chow/preamble.tex and
% sources/stacks-etale-cohomology/preamble.tex.  The block below is copied mechanically, line for line,
% from sources/stacks-chow/preamble.tex lines 14-18 (the Xy-pic setup) and 154-217 (the native theorem
% declarations together with the mathematical macros).  Nothing was added, removed or redefined; the
% wrapper class and the amsmath/amssymb/amsthm/graphicx base packages are supplied by the assembler.
% Chapter: more-algebra.tex -- More on Algebra.
% For commutative diagrams we use Xy-pic
\usepackage[all]{xy}

% We use 2cell for 2-commutative diagrams.
\xyoption{2cell}

% Theorem environments.
%
\theoremstyle{plain}
\newtheorem{theorem}[subsection]{Theorem}
\newtheorem{proposition}[subsection]{Proposition}
\newtheorem{lemma}[subsection]{Lemma}

\theoremstyle{definition}
\newtheorem{definition}[subsection]{Definition}
\newtheorem{example}[subsection]{Example}
\newtheorem{exercise}[subsection]{Exercise}
\newtheorem{situation}[subsection]{Situation}

\theoremstyle{remark}
\newtheorem{remark}[subsection]{Remark}
\newtheorem{remarks}[subsection]{Remarks}

\numberwithin{equation}{subsection}

% Macros
%
\def\lim{\mathop{\mathrm{lim}}\nolimits}
\def\colim{\mathop{\mathrm{colim}}\nolimits}
\def\Spec{\mathop{\mathrm{Spec}}}
\def\Hom{\mathop{\mathrm{Hom}}\nolimits}
\def\Ext{\mathop{\mathrm{Ext}}\nolimits}
\def\SheafHom{\mathop{\mathcal{H}\!\mathit{om}}\nolimits}
\def\SheafExt{\mathop{\mathcal{E}\!\mathit{xt}}\nolimits}
\def\Sch{\mathit{Sch}}
\def\Mor{\mathop{\mathrm{Mor}}\nolimits}
\def\Ob{\mathop{\mathrm{Ob}}\nolimits}
\def\Sh{\mathop{\mathit{Sh}}\nolimits}
\def\NL{\mathop{N\!L}\nolimits}
\def\CH{\mathop{\mathrm{CH}}\nolimits}
\def\proetale{{pro\text{-}\acute{e}tale}}
\def\etale{{\acute{e}tale}}
\def\QCoh{\mathit{QCoh}}
\def\Ker{\mathop{\mathrm{Ker}}}
\def\Im{\mathop{\mathrm{Im}}}
\def\Coker{\mathop{\mathrm{Coker}}}
\def\Coim{\mathop{\mathrm{Coim}}}

% Boxtimes
%
\DeclareMathSymbol{\boxtimes}{\mathbin}{AMSa}{"02}

%
% Macros for moduli stacks/spaces
%
\def\QCohstack{\mathcal{QC}\!\mathit{oh}}
\def\Cohstack{\mathcal{C}\!\mathit{oh}}
\def\Spacesstack{\mathcal{S}\!\mathit{paces}}
\def\Quotfunctor{\mathrm{Quot}}
\def\Hilbfunctor{\mathrm{Hilb}}
\def\Curvesstack{\mathcal{C}\!\mathit{urves}}
\def\Polarizedstack{\mathcal{P}\!\mathit{olarized}}
\def\Complexesstack{\mathcal{C}\!\mathit{omplexes}}
% \Pic is the operator that assigns to X its picard group, usage \Pic(X)
% \Picardstack_{X/B} denotes the Picard stack of X over B
% \Picardfunctor_{X/B} denotes the Picard functor of X over B
\def\Pic{\mathop{\mathrm{Pic}}\nolimits}
\def\Picardstack{\mathcal{P}\!\mathit{ic}}
\def\Picardfunctor{\mathrm{Pic}}
\def\Deformationcategory{\mathcal{D}\!\mathit{ef}}
\usepackage{cleveref}
\crefname{originallemma}{original lemma}{original lemmas}
\Crefname{originallemma}{Original lemma}{Original lemmas}
\makeatletter
\newlabel{more-algebra-lemma-characterize-injective-bis}{{0AVF}{source}{}{sourcepointer-1}{}}
\newlabel{more-algebra-lemma-characterize-injective-bis@cref}{{[originallemma][0][]0AVF}{[1][0][]source}{}{}{}}
\makeatother
\crefname{originallemma}{original lemma}{original lemmas}
\Crefname{originallemma}{Original lemma}{Original lemmas}
\makeatletter
\newlabel{algebra-lemma-reverse-long-exact-seq-ext}{{065P}{source}{}{sourcepointer-2}{}}
\newlabel{algebra-lemma-reverse-long-exact-seq-ext@cref}{{[originallemma][0][]065P}{[1][0][]source}{}{}{}}
\makeatother
\crefname{originallemma}{original lemma}{original lemmas}
\Crefname{originallemma}{Original lemma}{Original lemmas}
\makeatletter
\newlabel{algebra-lemma-characterize-projective}{{05CF}{source}{}{sourcepointer-3}{}}
\newlabel{algebra-lemma-characterize-projective@cref}{{[originallemma][0][]05CF}{[1][0][]source}{}{}{}}
\makeatother
\title{A second characterisation of injective modules through Baer's criterion}
\author{The Stacks Project Authors}
\date{stacks-project:more-algebra.tex:local-snapshot-20261012; GNU FDL 1.2 or later}
\begin{document}
\maketitle
\setcounter{subsection}{1}
\noindent\textit{Editorial note.} All delivered mathematical text is the author's own, copied from the pinned Stacks Project chapter file more-algebra.tex; no statement or proof step was rewritten, reordered or removed.\par\medskip
\noindent\textit{Editorial note.} The chapter file pinned-original-source.tex carries only \textbackslash{}input\{preamble\} on its first line, so no preamble text lives inside the chapter. The retained file sources/stacks-more-algebra/retained-preamble.tex is a mechanical, line-for-line copy of the project's single shared preamble as staged in the snapshot (Xy-pic setup plus theorem declarations and mathematical macros). Nothing was added, removed or redefined, and the excerpt serves only as preamble, never as a span coordinate.\par\medskip
\noindent\textit{Editorial note.} Both counted spans start at the \textbackslash{}begin line of their environment, so the statement and the author's proof are delivered whole; surrounding prose of the section is deliberately outside the counted spans and is carried as standard prerequisites instead.\par\medskip
\noindent\textit{Editorial note.} The citation fields (source\_id, source\_version, authors, source\_url, license.name, source\_title) are identical to those of every other card this lane has produced from this chapter, so the batch selector treats them as one source.\par\medskip
\noindent\textit{Allowed standard prerequisites (editorial).}\begin{itemize}
\item Injective module and the extension property for inclusions of submodules. defined earlier in the chapter, source lines 13552-13690, before the counted statement. Supplies the property being characterised.
\item Baer's criterion: it suffices to test the extension property on ideals of the ring. recalled and used at proof lines 13703-13720. Reduces the verification to the case of ideals.
\item Zorn's lemma applied to partial extensions and the reduction to a cyclic extension. used at proof lines 13720-13750. Produces the maximal counterexample and the single step that has to be verified.
\end{itemize}\medskip
\ifdefined\hypertarget\hypertarget{sourcepointer-1}{}\fi
\noindent\textit{Source reference.} Original lemma 0AVF: focused-seeds/AGP/sources/stacks-snapshot-20261011/tags.txt, source lines 13683--13683. The lemma's own canonical tag, quoted so that readers can resolve the label inside the delivered text; it points at the counted statement itself and carries no proof content.\par\medskip
\ifdefined\hypertarget\hypertarget{sourcepointer-2}{}\fi
\noindent\textit{Source reference.} Original lemma 065P: focused-seeds/AGP/sources/stacks-snapshot-20261011/tags.txt, source lines 7908--7908. Navigation pointer to a statement the counted proof refers to but does not contain; nothing load-bearing is delegated to it.\par\medskip
\ifdefined\hypertarget\hypertarget{sourcepointer-3}{}\fi
\noindent\textit{Source reference.} Original lemma 05CF: focused-seeds/AGP/sources/stacks-snapshot-20261011/tags.txt, source lines 6918--6918. Navigation pointer to a statement the counted proof refers to but does not contain; nothing load-bearing is delegated to it.\par\medskip
\begin{lemma}
\label{lemma-characterize-injective-bis}
Let $R$ be a ring. Let $J$ be an $R$-module.
The following are equivalent
\begin{enumerate}
\item $J$ is injective,
\item $\Ext^1_R(R/I, J) = 0$ for every ideal $I \subset R$, and
\item for an ideal $I \subset R$ and module map $I \to J$
there exists an extension $R \to J$.
\end{enumerate}
\end{lemma}

\begin{proof}
If $I \subset R$ is an ideal, then the short exact sequence
$0 \to I \to R \to R/I \to 0$ gives an exact sequence
$$
\Hom_R(R, J) \to
\Hom_R(I, J) \to
\Ext^1_R(R/I, J) \to 0
$$
by Algebra, Lemma \ref{algebra-lemma-reverse-long-exact-seq-ext}
and the fact that $\Ext^1_R(R, J) = 0$ as $R$ is projective
(Algebra, Lemma \ref{algebra-lemma-characterize-projective}).
Thus (2) and (3) are equivalent. In this proof
we will show that (1) $\Leftrightarrow$ (3) which is known
as Baer's criterion.

\medskip\noindent
Assume (1). Given a module map $I \to J$ as in (3) we find
the extension $R \to J$ because the map
$\Hom_R(R, J) \to \Hom_R(I, J)$ is surjective
by definition.

\medskip\noindent
Assume (3). Let $M \subset N$ be an inclusion of $R$-modules.
Let $\varphi : M \to J$ be a homomorphism. We will show that $\varphi$
extends to $N$ which finishes the proof of the lemma.
Consider the set of homomorphisms $\varphi' : M' \to J$
with $M \subset M' \subset N$ and $\varphi'|_M = \varphi$.
Define $(M', \varphi') \geq (M'', \varphi'')$ if
and only if $M' \supset M''$ and $\varphi'|_{M''} = \varphi''$.
If $(M_i, \varphi_i)_{i \in I}$ is a totally
ordered collection of such pairs, then we obtain a map
$\bigcup_{i \in I} M_i \to J$ defined by $a \in M_i$
maps to $\varphi_i(a)$. Thus Zorn's lemma applies.
To conclude we have to show that if the pair
$(M', \varphi')$ is maximal then $M' = N$.
In other words, it suffices to show, given
any subgroup $M \subset N$, $M \not = N$ and
any $\varphi : M \to J$, then we can find
$\varphi' : M' \to J$ with $M \subset M' \subset N$
such that (a) the inclusion $M \subset M'$ is strict, and
(b) the morphism $\varphi'$ extends $\varphi$.

\medskip\noindent
To prove this, pick $x \in N$, $x \not \in M$.
Let $I = \{f \in R \mid fx \in M\}$. This is an ideal of $R$.
Define a homomorphism $\psi : I \to J$ by $f \mapsto \varphi(fx)$.
Extend to a map $\tilde\psi : R \to J$ which is possible by assumption (3).
By our choice of $I$ the kernel of
$M \oplus R \to J$, $(y, f) \mapsto \varphi(y) + \tilde\psi(f)$
contains the kernel of the map $M \oplus R \to N$,
$(y, f) \mapsto y + fx$. Hence this homomorphism factors
through the image $M' = M + Rx$ and this extends the given homomorphism
as desired.
\end{proof}

\end{document}
